% Oversells «Diniy matn» shabloni: arabcha matn oʻngdan chapga (LuaLaTeX tavsiya etiladi, XeLaTeX ham ishlaydi).
% Har bir oyat — chap ustunda tarjima, oʻng ustunda arabcha matn (paracol).
% \oyat{arabcha matn}{raqam}{tarjima} — bitta oyat; uslublar sahifa.sty da
% («Sahifa» tugmasi). Arabcha soʻz matn ichida: \arabcha{...}.
\documentclass[11pt,twoside]{article}
\usepackage[paperwidth=148mm,paperheight=210mm,headheight=14pt]{geometry}
\usepackage{iftex}
% Right to left: babel's bidi=basic in LuaLaTeX (best), XeTeX's own in XeLaTeX.
\ifLuaTeX\PassOptionsToPackage{bidi=basic}{babel}\else\PassOptionsToPackage{bidi=default}{babel}\fi
\usepackage{babel}
\babelprovide[import=uz-Latn, main]{uzbek}
\babelprovide[import, mapdigits]{arabic}
\babelfont{rm}{Noto Serif}
\ifLuaTeX\babelfont[arabic]{rm}[Renderer=HarfBuzz]{Amiri}\else\babelfont[arabic]{rm}{Amiri}\fi
\def\uzbekhyphenmins{22}
\usepackage{paracol}
\usepackage{sahifa} % Oversells sahifa dizayni (sahifa.sty)

\newcommand\arabcha[1]{\foreignlanguage{arabic}{#1}}
% One verse: the translation on the left, the Arabic text on the right, both on one line.
\newcommand\oyat[3]{%
  \begin{uslub}{tarjima}\textbf{\matnuslubi{raqam}{#2.}}~#3\end{uslub}%
  \switchcolumn
  \begin{otherlanguage}{arabic}\begin{uslub}{oyat}#1~\matnuslubi{raqam}{﴿\localenumeral{digits}{#2}﴾}\end{uslub}\end{otherlanguage}%
  \switchcolumn*}

\begin{document}

\section*{Fotiha surasi}
\markboth{Fotiha surasi}{}

\begin{otherlanguage}{arabic}
\begin{uslub}{basmala}
بِسْمِ ٱللَّهِ ٱلرَّحْمَٰنِ ٱلرَّحِيمِ
\end{uslub}
\end{otherlanguage}

\columnratio{0.45}
\begin{paracol}{2}
\oyat{بِسْمِ ٱللَّهِ ٱلرَّحْمَٰنِ ٱلرَّحِيمِ}{1}{Mehribon va rahmli Alloh nomi bilan.}
\oyat{ٱلْحَمْدُ لِلَّهِ رَبِّ ٱلْعَٰلَمِينَ}{2}{Hamd olamlarning Rabbi Allohgadir.}
\oyat{ٱلرَّحْمَٰنِ ٱلرَّحِيمِ}{3}{U mehribon va rahmlidir.}
\oyat{مَٰلِكِ يَوْمِ ٱلدِّينِ}{4}{Jazo kunining egasidir.}
\oyat{إِيَّاكَ نَعْبُدُ وَإِيَّاكَ نَسْتَعِينُ}{5}{Faqat Senga ibodat qilamiz va faqat Sendan yordam soʻraymiz.}
\oyat{ٱهْدِنَا ٱلصِّرَٰطَ ٱلْمُسْتَقِيمَ}{6}{Bizni toʻgʻri yoʻlga boshlagin,}
\oyat{صِرَٰطَ ٱلَّذِينَ أَنْعَمْتَ عَلَيْهِمْ غَيْرِ ٱلْمَغْضُوبِ عَلَيْهِمْ وَلَا ٱلضَّآلِّينَ}{7}{Oʻzing neʼmat bergan zotlarning yoʻliga; gʻazabga uchraganlarning va adashganlarning yoʻliga emas.}
\end{paracol}

\begin{uslub}{izoh}
Tarjima maʼnoning qisqa bayoni. Arabcha soʻzni matn ichida ham yozish mumkin:
\arabcha{ٱلْحَمْدُ لِلَّهِ} — «hamd Allohgadir».
\end{uslub}

\end{document}
